\documentclass[12pt]{article}
\usepackage[utf8]{inputenc}
\usepackage[T1]{fontenc}
\usepackage{hyperref}
\title{Nganakarti Language Resources for K-12 Curriculum Development}
\author{Compiled by Claude Sonnet 5.5 (claude-sonnet-5-5) for Andr\'as Kornai}
\date{}
\begin{document}
\maketitle

\section{Introduction}
Nganakarti (ISO 639-3 xnk) is an extinct Aboriginal Australian language. Searches found no page under the spelling ``Nganakarti''. The code corresponds to the language Wikipedia calls Nhanhagardi (also Nanakarti and similar variants), from the Champion Bay area of Western Australia. No speakers have been recorded since 1975. Its classification has been questioned: one source treats it as a dialect of Njungar, another as a name for Widi. Very little documentation was found, so a curriculum cannot be built from the language itself at present.

\section{Primary Educational Resources}
No primers, textbooks or school programmes were found in search results.

\section{Linguistic and Reference Materials}
\begin{itemize}
\item Nhanhagardi language (Wikipedia overview): \url{https://en.wikipedia.org/wiki/Nhanhagardi_language}
\item AIATSIS Austlang entry A93, Nhanhagardi, with the language's status and alternative names: \url{https://collection.aiatsis.gov.au/austlang/language/A93}
\end{itemize}

\section{Religious and Community Resources}
None found.

\section{Digital and Audio Resources}
The search summary says AIATSIS holds a short recording of a song in Nhanhagardi. I did not see a direct link to it, so it is listed only through the Austlang entry above.

\section{Curriculum Development Context}
With no known speakers and almost no documentation, a K-12 programme could realistically only present the language's history and place, using the AIATSIS record. Any language content would need a request to AIATSIS and consultation with Aboriginal communities of the Geraldton region. Nothing found supports claiming a usable vocabulary or grammar.
\end{document}